\section{Summary and Conclusions}

In this study, we have analyzed the gamma-ray emission from the Milky Way's globular clusters, using 15.8 years of publicly available data collected by the Fermi Gamma-Ray Space Telescope. We report the robust detection of 56 globular clusters with a statistically significant test statistic of TS > 25, 8 of which are not contained in previous gamma-ray source catalogs. 

Making use of previously established correlations between the mean number of millisecond pulsars in a globular cluster and the stellar encounter rate or visible luminosity of that cluster, we perform a determination of the gamma-ray luminosity function of the millisecond pulsars in the Milky Way's globular cluster system. To this end, we have considered 87 globular clusters with high stellar encounter rates and/or visible luminosities, and combined this information with our measurements of their gamma-ray luminosities. We then fitted the data to a log-normal luminosity function, finding it to favor a mean pulsar luminosity of $\langle L_{\gamma}\rangle \sim (1-8)\times 10^{33}\, \textrm{erg/s}$ (integrated between 0.1 and 100 GeV), and a width of $\sigma_L \sim 1.4-2.7$. Our analysis also indicates that the Milky Way's globular cluster system contains a total of $\sim 400-1500$ gamma-ray emitting millisecond pulsars. The luminosity function favored by this analysis is consistent with that exhibited by millisecond pulsars found in the Galactic Plane~\cite{Holst:2024fvb}.  

The leading astrophysical interpretation of the Galactic Center Gamma-Ray Excess has long been that this signal could be generated by large population of centrally-located millisecond pulsars. Very few pulsars have been detected in the Inner Galaxy, however, forcing us to place strong constraints on the gamma-ray luminosity function of any pulsar population that could be responsible for the excess emission. In particular, such a pulsar population would have to be systematically less bright than those found in the Galactic Plane~\cite{Holst:2024fvb}. 

It has been suggested that the Inner Galaxy pulsar population could be older and therefore less luminous, on average, than pulsars found in the Galactic Plane. The stars in globular clusters are also very old, however. Thus if the Inner Galaxy pulsar population is very faint, one would expect the same to be true among the pulsars found in globular clusters. The results of this study indicate that this is not the case. This result poses a challenge to pulsar interpretations of the Galactic Center Gamma-Ray Excess and, by extension, potentially bolsters dark matter interpretations of this long-standing signal. 

\bigskip


\appendix

